\documentclass[10pt,a4paper]{article}
\usepackage[T1]{fontenc}
\usepackage[utf8]{inputenc}
\usepackage[brazilian]{babel}
\usepackage[a4paper,margin=1.5cm]{geometry}
\usepackage{booktabs}
\usepackage{xcolor}
\usepackage{microtype}
\usepackage[hidelinks]{hyperref}
\pagestyle{empty}
\setlength{\parindent}{0pt}
\setlength{\parskip}{2pt}
% Paleta monocromática: só a escala gray do xcolor, para impressão em toner.
\definecolor{cinza}{gray}{0.35}
\definecolor{cinzaclaro}{gray}{0.94}
\definecolor{borda}{gray}{0.75}
\definecolor{fundocard}{gray}{0.98}
\definecolor{barras}{gray}{0.15}
\newcommand{\barra}[1]{\textcolor{barras}{\rule{#1mm}{2.4pt}}}
% Cartão de indicador: rótulo, número em destaque e rodapé de contexto.
\newcommand{\kpicard}[3]{%
  \fcolorbox{borda}{fundocard}{%
    \begin{minipage}[c][1.4cm][c]{\dimexpr0.32\textwidth-2\fboxsep-2\fboxrule\relax}
      \centering
      {\scriptsize\bfseries\color{cinza} #1}\\[2pt]
      {\Large\bfseries #2}\\[1.5pt]
      {\tiny\color{cinza} #3}%
    \end{minipage}%
  }%
}
\begin{document}

{\Large\bfseries Pedro de Toledo --- climatologia das médias}\\[2pt]
{\footnotesize\color{cinza} Estação CIIAGRO/IAC ciiagro:300 \,\textbullet\, 99 meses (2018-08 a 2026-10)}\\[1pt]
{\footnotesize \textbf{Atenção à série curta:} apenas 7 anos com os 12 meses fechados --- referência do período medido, não padrão consolidado.}

\vspace{4pt}
\colorbox{cinzaclaro}{%
  \parbox{\dimexpr\textwidth-2\fboxsep}{%
    \footnotesize
    \textbf{Leia antes de usar:} os números abaixo são \textbf{médias de vários anos}, obtidas das medições mensais do CIIAGRO/IAC. \textbf{Não é previsão meteorológica} para os próximos dias nem para esta safra.\par\vspace{2pt}
    \textbf{Aplicação prática:} este é o comportamento médio esperado da estação ---
    base técnica para o calendário agrícola, janelas de calagem, o preparo do solo e
    o manejo conservacionista.%
  }%
}
\vspace{4pt}
\noindent
\kpicard{Chuva anual média}{1.240 mm/ano}{em todo o ano}\hfill
\kpicard{Mês mais chuvoso}{219 mm}{em janeiro}\hfill
\kpicard{Mês mais seco}{34 mm}{em julho}

\vspace{2pt}
\noindent
\kpicard{Temperatura média}{23,9 \textdegree{}C}{média dos 12 meses}\hfill
\kpicard{Mês mais quente}{27,9 \textdegree{}C}{em janeiro}\hfill
\kpicard{Mês mais frio}{20,2 \textdegree{}C}{em julho}

\vspace{4pt}
\begin{center}
{\bfseries\small Tabela 1: Chuva média por mês do ano}\\[2pt]
{\small
\begin{tabular}{@{}l r l r r r@{}}
\toprule
mês & chuva (mm) & perfil & t mín & t méd & t máx \\
\midrule
janeiro    &   219 & \barra{26.0}   & 22,5 & 27,9 & 33,4 \\
fevereiro  &   187 & \barra{22.2}   & 21,8 & 27,5 & 33,3 \\
março      &   163 & \barra{19.4}   & 22,0 & 27,2 & 32,4 \\
abril      &    78 & \barra{9.2}    & 19,3 & 24,7 & 30,1 \\
maio       &    90 & \barra{10.7}   & 15,9 & 21,9 & 28,0 \\
junho      &    63 & \barra{7.5}    & 15,5 & 21,1 & 26,7 \\
julho      &    34 & \barra{4.1}    & 14,2 & 20,2 & 26,2 \\
agosto     &    46 & \barra{5.5}    & 15,5 & 20,9 & 26,2 \\
setembro   &    56 & \barra{6.7}    & 17,6 & 22,5 & 27,4 \\
outubro    &    79 & \barra{9.4}    & 19,2 & 23,5 & 27,8 \\
novembro   &   121 & \barra{14.3}   & 19,6 & 24,6 & 29,6 \\
dezembro   &   129 & \barra{15.4}   & 18,1 & 25,1 & 32,6 \\
\midrule
ano & 1.240 & & & 23,9 & \\
\bottomrule
\end{tabular}}
\end{center}

\begin{center}
{\bfseries\small Tabela 2: Chuva por estação do ano}\\[2pt]
{\small
\begin{tabular}{@{}l l r r r l l@{}}
\toprule
estação & meses & total (mm) & \% do ano & média/mês & maior mês & menor mês \\
\midrule
Verão      & dez--fev &   535 &  43\% &  178 & 219 (jan) & 129 (dez) \\
Outono     & mar--mai &   331 &  27\% &  110 & 163 (mar) & 78 (abr) \\
Inverno    & jun--ago &   143 &  12\% &   48 & 63 (jun) & 34 (jul) \\
Primavera  & set--nov &   256 &  21\% &   85 & 121 (nov) & 56 (set) \\
\bottomrule
\end{tabular}}
\end{center}

\vspace{1pt}
{\footnotesize\textbf{Antes dos solstícios} (acumulado desde janeiro):
inverno: janeiro a maio somam 736 mm (59\% do ano) \,\textbullet\, verão: junho a novembro somam 400 mm (32\% do ano). Dezembro fica depois do solstício de verão e entra só no total
anual.}

\vspace{2pt}
{\bfseries\small Como usar no planejamento nesta estação}

\footnotesize\begin{itemize}\setlength{\itemsep}{0.5pt}
  \item \textbf{Amostragem de solo e preparo mecânico}: Em abril, maio, junho, julho, agosto, setembro e outubro a chuva média fica entre 34 e 90 mm/mês \textemdash{} é a janela de menor encharcamento para retirar terra para laudo, arar e incorporar calcário com antecedência.
  \item \textbf{Calagem em cobertura (pasto e perene já formado)}: A partir de novembro, quando a média volta a passar de 120 mm/mês: aí há água para o corretivo descer no perfil sem enxurrada.
  \item \textbf{Plantio e transplante}: O verão (dezembro a fevereiro) concentra 535 mm, 43\% da chuva do ano \textemdash{} favorece o pegamento de mudas sem irrigação pesada; em áreas declivosas, reforce a conservação de solo nesse período.
  \item \textbf{Excesso de chuva e conservação}: janeiro passa de 200 mm em média (219 mm): canteiros altos, drenagem limpa e cobertura morta antes da temporada, para não perder solo e adubo no escoamento.
  \item \textbf{Calor, frio e fitossanidade}: janeiro tem a maior média (27,9 \textdegree{}C) e julho a menor (20,2 \textdegree{}C, com mínima média de 14,2 \textdegree{}C). Antecipe o monitoramento de doenças fúngicas nos meses quentes e úmidos e proteja culturas sensíveis ao frio nos mais frios.
\end{itemize}

\vspace{1pt}
{\bfseries\small Vale para qualquer município da região}

\scriptsize\begin{itemize}\setlength{\itemsep}{0.5pt}
  \item \textbf{Média não é previsão}: Descrevem o padrão de vários anos, não o tempo deste ano: use para escolher a janela de cada operação, e decida no campo.
  \item \textbf{Corrija antes de precisar}: Calcário reage em meses; aplicado na janela seca, chega ao plantio já reagido.
  \item \textbf{Chuva em mm não é água no solo}: Textura, declive e cobertura morta decidem quanto do que cai chega à raiz.
  \item \textbf{Série curta é pista, não padrão}: Poucos anos fechados oscilam muito: compare com a estação vizinha antes de fixar calendário.
\end{itemize}

\vspace{1pt}
{\tiny\color{cinza}
Fonte: CIIAGRO/IAC (\href{http://ciiagro.org.br}{ciiagro.org.br}), série mensal;
chuva é o total do mês. Painel:
\href{https://climalitoralsul.netlify.app/\#pedro-de-toledo}%
{\texttt{climalitoralsul.netlify.app/\#pedro-de-toledo}}.
Umidade relativa omitida: a fonte usa \texttt{0} para ``sem sensor''.  6 leitura(s) fora da faixa de sanidade excluída(s) da média (tmed 3, tmax 3); o banco as preserva.
médias do período 2018\textendash{}2026 (8.2 anos de dado).
Gerado a partir de \texttt{dados.json} de 2026-10-30.}

\end{document}
